\documentclass[10pt,a4paper]{amsart}
\usepackage{amsmath,amssymb,amsthm,graphicx}
% For commutative diagrams we use Xy-pic
\usepackage[all]{xy}

% We use 2cell for 2-commutative diagrams.
\xyoption{2cell}
% Theorem environments.
%
\theoremstyle{plain}
\newtheorem{theorem}[subsection]{Theorem}
\newtheorem{proposition}[subsection]{Proposition}
\newtheorem{lemma}[subsection]{Lemma}

\theoremstyle{definition}
\newtheorem{definition}[subsection]{Definition}
\newtheorem{example}[subsection]{Example}
\newtheorem{exercise}[subsection]{Exercise}
\newtheorem{situation}[subsection]{Situation}

\theoremstyle{remark}
\newtheorem{remark}[subsection]{Remark}
\newtheorem{remarks}[subsection]{Remarks}

\numberwithin{equation}{subsection}

% Macros
%
\def\lim{\mathop{\mathrm{lim}}\nolimits}
\def\colim{\mathop{\mathrm{colim}}\nolimits}
\def\Spec{\mathop{\mathrm{Spec}}}
\def\Hom{\mathop{\mathrm{Hom}}\nolimits}
\def\Ext{\mathop{\mathrm{Ext}}\nolimits}
\def\SheafHom{\mathop{\mathcal{H}\!\mathit{om}}\nolimits}
\def\SheafExt{\mathop{\mathcal{E}\!\mathit{xt}}\nolimits}
\def\Sch{\mathit{Sch}}
\def\Mor{\mathop{\mathrm{Mor}}\nolimits}
\def\Ob{\mathop{\mathrm{Ob}}\nolimits}
\def\Sh{\mathop{\mathit{Sh}}\nolimits}
\def\NL{\mathop{N\!L}\nolimits}
\def\CH{\mathop{\mathrm{CH}}\nolimits}
\def\proetale{{pro\text{-}\acute{e}tale}}
\def\etale{{\acute{e}tale}}
\def\QCoh{\mathit{QCoh}}
\def\Ker{\mathop{\mathrm{Ker}}}
\def\Im{\mathop{\mathrm{Im}}}
\def\Coker{\mathop{\mathrm{Coker}}}
\def\Coim{\mathop{\mathrm{Coim}}}

% Boxtimes
%
\DeclareMathSymbol{\boxtimes}{\mathbin}{AMSa}{"02}

%
% Macros for moduli stacks/spaces
%
\def\QCohstack{\mathcal{QC}\!\mathit{oh}}
\def\Cohstack{\mathcal{C}\!\mathit{oh}}
\def\Spacesstack{\mathcal{S}\!\mathit{paces}}
\def\Quotfunctor{\mathrm{Quot}}
\def\Hilbfunctor{\mathrm{Hilb}}
\def\Curvesstack{\mathcal{C}\!\mathit{urves}}
\def\Polarizedstack{\mathcal{P}\!\mathit{olarized}}
\def\Complexesstack{\mathcal{C}\!\mathit{omplexes}}
% \Pic is the operator that assigns to X its picard group, usage \Pic(X)
% \Picardstack_{X/B} denotes the Picard stack of X over B
% \Picardfunctor_{X/B} denotes the Picard functor of X over B
\def\Pic{\mathop{\mathrm{Pic}}\nolimits}
\def\Picardstack{\mathcal{P}\!\mathit{ic}}
\def\Picardfunctor{\mathrm{Pic}}
\def\Deformationcategory{\mathcal{D}\!\mathit{ef}}
\usepackage{cleveref}
\crefname{originalreference}{original reference}{original references}
\Crefname{originalreference}{Original reference}{Original references}
\makeatletter
\newlabel{definition-locally-acyclic}{{0GJP}{source}{}{sourcepointer-1}{}}
\newlabel{definition-locally-acyclic@cref}{{[originalreference][0][]0GJP}{[1][0][]source}{}{}{}}
\makeatother
\crefname{originalreference}{original reference}{original references}
\Crefname{originalreference}{Original reference}{Original references}
\makeatletter
\newlabel{lemma-base-change-f-star-general-stalks}{{0EZU}{source}{}{sourcepointer-2}{}}
\newlabel{lemma-base-change-f-star-general-stalks@cref}{{[originalreference][0][]0EZU}{[1][0][]source}{}{}{}}
\makeatother
\crefname{originalreference}{original reference}{original references}
\Crefname{originalreference}{Original reference}{Original references}
\makeatletter
\newlabel{theorem-smooth-base-change}{{0EYU}{source}{}{sourcepointer-3}{}}
\newlabel{theorem-smooth-base-change@cref}{{[originalreference][0][]0EYU}{[1][0][]source}{}{}{}}
\makeatother
\title{Smooth morphisms are universally locally acyclic}
\author{The Stacks Project Authors}
\date{stacks-project:etale-cohomology.tex:local-snapshot-20261011; GNU FDL 1.2 or later}
\begin{document}
\maketitle
\noindent\textit{Editorial note.} All delivered mathematical text is the author's own, copied from the pinned Stacks Project chapter file etale-cohomology.tex; no statement or proof step was rewritten, reordered or removed.\par\medskip
\noindent\textit{Editorial note.} The counted proof is the author's complete argument: the proof is written out in full: the reduction through the definition of local acyclicity, the application of smooth base change, the displayed chain of cohomology identifications, and the final vanishing computation.\par\medskip
\noindent\textit{Editorial note.} The label proposition-smooth-locally-acyclic names the counted statement, and the proof environment immediately following it in the source is delivered verbatim as its proof.\par\medskip
\noindent\textit{Editorial note.} The counted statement is carried by the author's own proposition environment in the chapter source.\par\medskip
\noindent\textit{Editorial note.} The statement is an author Proposition environment, delivered verbatim; the three cited inputs are kept as external pointers with their canonical Stacks tags.\par\medskip
\noindent\textit{Allowed standard prerequisites (editorial).}\begin{itemize}
\item Local acyclicity as a property of a ring map. Same chapter, the definition of local acyclicity and the discussion after it, which reduces the property to a statement about the map of strict henselisations. Supplies the formulation used to reduce the statement to the strictly henselian case and the target vanishing.
\item Smooth base change. Same chapter, the smooth base change theorem for torsion coefficients invertible on the base. Supplies the isomorphism f\textasciicircum{}\{-1\}R\textasciicircum{}qg\_*F = R\textasciicircum{}qh\_*e\textasciicircum{}\{-1\}F for the constant coefficients Z/nZ.
\item The stalk form of the base change statement. Same chapter, the lemma reformulating the base change isomorphism as an equality of cohomology groups over the strict henselisations. Supplies the middle equality of the displayed chain.
\end{itemize}\medskip
\ifdefined\hypertarget\hypertarget{sourcepointer-1}{}\fi
\noindent\textit{Source reference.} Original reference 0GJP: sources/stacks-etale-cohomology/pinned-original-source.tex, source lines 19110--19110. Prerequisite pointer invoked at proof line(s) 19181; it names a standard fact of the same chapter that the counted argument uses, and every step of the counted proof built on it is retained verbatim in the delivered body.\par\medskip
\ifdefined\hypertarget\hypertarget{sourcepointer-2}{}\fi
\noindent\textit{Source reference.} Original reference 0EZU: sources/stacks-etale-cohomology/pinned-original-source.tex, source lines 16497--16497. Prerequisite pointer invoked at proof line(s) 19190; it names a standard fact of the same chapter that the counted argument uses, and every step of the counted proof built on it is retained verbatim in the delivered body.\par\medskip
\ifdefined\hypertarget\hypertarget{sourcepointer-3}{}\fi
\noindent\textit{Source reference.} Original reference 0EYU: sources/stacks-etale-cohomology/pinned-original-source.tex, source lines 17446--17446. Prerequisite pointer invoked at proof line(s) 19203; it names a standard fact of the same chapter that the counted argument uses, and every step of the counted proof built on it is retained verbatim in the delivered body.\par\medskip
\begin{proposition}
\label{proposition-smooth-locally-acyclic}
Let $f : X \to S$ be a smooth morphism of schemes.
Then $f$ is universally locally acyclic.
\end{proposition}

\begin{proof}
Since the base change of a smooth morphism is smooth, it suffices
to show that smooth morphisms are locally acyclic.
Let $\overline{x}$ be a geometric point of $X$ with image
$\overline{s} = f(\overline{x})$. Let
$\overline{t}$ be a geometric point of
$\Spec(\mathcal{O}^{sh}_{S, f(\overline{x})})$.
Since we are trying to prove a property of the ring map
$\mathcal{O}^{sh}_{S, \overline{s}} \to \mathcal{O}^{sh}_{X, \overline{x}}$
(see discussion following Definition \ref{definition-locally-acyclic})
we may and do replace $f : X \to S$ by the base change
$X \times_S \Spec(\mathcal{O}^{sh}_{S, \overline{s}}) \to
\Spec(\mathcal{O}^{sh}_{S, \overline{s}})$.
Thus we may and do assume that $S$ is the spectrum of a strictly
henselian local ring and that $\overline{s}$ lies over
the closed point of $S$.

\medskip\noindent
We will apply Lemma \ref{lemma-base-change-f-star-general-stalks}
to the diagram
$$
\xymatrix{
X \ar[d]_f & X_{\overline{t}} \ar[l]^h \ar[d]^e \\
S & \overline{t} \ar[l]_g
}
$$
and the sheaf $\mathcal{F} = \underline{M}$ where $M = \mathbf{Z}/n\mathbf{Z}$
for some integer $n$ prime to the characteristic of the residue field of
$\overline{x}$. We know that the map
$f^{-1}R^qg_*\mathcal{F} \to R^qh_*e^{-1}\mathcal{F}$
is an isomorphism by smooth base change, see
Theorem \ref{theorem-smooth-base-change} (the assumption
on torsion holds by our choice of $n$).
Thus Lemma \ref{lemma-base-change-f-star-general-stalks}
gives us the middle equality in
$$
H^q(F_{\overline{x}, \overline{t}}, \underline{M}) =
H^q(\Spec(\mathcal{O}^{sh}_{X, \overline{x}}) \times_S \overline{t},
\underline{M})
=
H^q(\Spec(\mathcal{O}^{sh}_{S, \overline{s}}) \times_S \overline{t},
\underline{M}) =
H^q(\overline{t}, \underline{M})
$$
For the outer two equalities we use that
$S = \Spec(\mathcal{O}^{sh}_{S, \overline{s}})$.
Since $\overline{t}$ is the spectrum of a separably
closed field we conclude that
$$
H^q(F_{\overline{x}, \overline{t}}, \underline{M}) =
\left\{
\begin{matrix}
M & \text{if} & q = 0 \\
0 & \text{if} & q \not = 0
\end{matrix}
\right.
$$
which is what we had to show (see discussion following
Definition \ref{definition-locally-acyclic}).
\end{proof}

\end{document}
